% ============================================================================
% ARBEITSBLATT DS 1 - Gruppe B: Schulsystem & Berufe (S. 102-103)
% ============================================================================

\documentclass[11pt,a4paper]{article}

\usepackage[utf8]{inputenc}
\usepackage[T1]{fontenc}
\usepackage[ngerman]{babel}
\usepackage[left=2cm,right=2cm,top=1.8cm,bottom=1.5cm]{geometry}
\usepackage{helvet}
\usepackage[table]{xcolor}
\usepackage{tabularx}
\usepackage{array}
\usepackage{enumitem}
\usepackage{tcolorbox}
\usepackage{fancyhdr}
\usepackage{lastpage}
\usepackage{amssymb}

\renewcommand{\familydefault}{\sfdefault}

\definecolor{spanisch}{HTML}{c41e3a}
\definecolor{grau}{HTML}{666666}
\definecolor{hellgrau}{HTML}{f5f5f5}
\definecolor{gruppeB}{HTML}{1565c0}

\pagestyle{fancy}
\fancyhf{}
\fancyhead[L]{\small\textcolor{grau}{Spanisch Spätbeginner -- Gruppe B}}
\fancyhead[R]{\small\textcolor{grau}{AB-01-B}}
\fancyfoot[C]{\small\thepage/\pageref{LastPage}}
\renewcommand{\headrulewidth}{0.5pt}

\newcommand{\luecke}[1]{\underline{\hspace{#1}}}
\newcommand{\punktebox}[1]{\hfill\fbox{\small \_\_/#1}}
\newcommand{\es}[1]{\textit{#1}}

\newtcolorbox{merkkasten}[1][]{
    colback=white,
    colframe=gruppeB,
    fonttitle=\bfseries\color{gruppeB},
    title=#1,
    boxrule=1.5pt,
    arc=0pt,
    left=5pt, right=5pt, top=5pt, bottom=5pt
}

\newtcolorbox{buchverweis}{
    colback=hellgrau,
    colframe=grau,
    boxrule=0.5pt,
    arc=0pt,
    left=5pt, right=5pt, top=3pt, bottom=3pt
}

\newcounter{aufgabe}
\newcommand{\aufgabe}[2]{%
    \stepcounter{aufgabe}%
    \vspace{0.8em}%
    \noindent\textbf{\theaufgabe.} #1 \punktebox{#2}%
    \vspace{0.3em}%
}

\begin{document}

% ============================================================================
% KOPFBEREICH
% ============================================================================
\noindent
\begin{tabularx}{\textwidth}{@{}Xr@{}}
    {\Large\textbf{\textcolor{gruppeB}{El sistema educativo}}} & Name: \luecke{5cm} \\[0.3em]
    {\small DS 1 -- Bucharbeit (S. 102--103)} & Datum: \luecke{3cm} \\
\end{tabularx}

\vspace{0.3em}
\hrule
\vspace{0.8em}

% ============================================================================
% MERKKASTEN 1: Konjunktionen
% ============================================================================
\begin{merkkasten}[Kasten 1: Kausale Konjunktionen]
\vspace{0.2em}

\begin{tabularx}{\textwidth}{@{}|p{3cm}|X|@{}}
\hline
\rowcolor{hellgrau}
\textbf{Konjunktion} & \textbf{Beispiel} \\
\hline
\textbf{ya que} (da, weil) & Ya que el sistema es diferente, hay que explicarlo. \\[0.4em]
\hline
\textbf{porque} (weil) & Estudio español porque \luecke{4cm}. \\[0.4em]
\hline
\textbf{mientras (que)} (während) & Mientras en Alemania hay varios tipos de escuela, en España... \\[0.4em]
\hline
\textbf{además} (außerdem) & Además, los estudiantes tienen que... \\[0.4em]
\hline
\end{tabularx}

\vspace{0.3em}
\end{merkkasten}

\vspace{0.8em}

% ============================================================================
% MERKKASTEN 2: Schulsystem Spanien
% ============================================================================
\begin{merkkasten}[Kasten 2: Sistema educativo español]
\vspace{0.2em}

\begin{tabularx}{\textwidth}{@{}|p{4cm}|c|X|@{}}
\hline
\rowcolor{hellgrau}
\textbf{Stufe} & \textbf{Alter} & \textbf{Beschreibung} \\
\hline
Educación Primaria & 6--12 & Grundschule (6 Jahre) \\[0.3em]
\hline
ESO (Educación Secundaria Obligatoria) & 12--16 & \luecke{4cm} \\[0.3em]
\hline
Bachillerato & 16--18 & \luecke{4cm} \\[0.3em]
\hline
FP (Formación Profesional) & 16+ & \luecke{4cm} \\[0.3em]
\hline
\end{tabularx}

\vspace{0.3em}
\textbf{Notensystem:} 0--10 Punkte (aprobado ab 5, sobresaliente = 9--10)
\vspace{0.2em}
\end{merkkasten}

\vspace{1em}

% ============================================================================
% AUFGABEN
% ============================================================================

\aufgabe{Erkläre das spanische Schulsystem. (Buch S. 102, Aufg. 6a)}{6}

\begin{buchverweis}
\textbf{Schreibe auf Spanisch.} Erkläre einem deutschen Austauschschüler, wie das spanische Schulsystem funktioniert.
\end{buchverweis}

\vspace{0.5em}
\begin{tabularx}{\textwidth}{@{}|X|@{}}
\hline
\\[0.5em]
\luecke{14cm} \\[0.7em]
\luecke{14cm} \\[0.7em]
\luecke{14cm} \\[0.7em]
\luecke{14cm} \\[0.5em]
\hline
\end{tabularx}

\vspace{0.8em}

\aufgabe{Vergleiche: Spanien vs. Deutschland. (Buch S. 102, Aufg. 6b)}{8}

\begin{buchverweis}
\textbf{Nutze die Konjunktionen} \es{mientras (que)}, \es{ya que}, \es{porque}, \es{además}.
\end{buchverweis}

\vspace{0.5em}
\begin{tabularx}{\textwidth}{@{}|p{6.5cm}|p{6.5cm}|@{}}
\hline
\rowcolor{hellgrau}
\textbf{España} & \textbf{Alemania} \\
\hline
Sistema centralizado & Sistema federal (Bundesländer) \\[0.3em]
\hline
Notas: 0--10 & Notas: \luecke{3cm} \\[0.3em]
\hline
ESO obligatoria hasta los 16 & \luecke{4cm} \\[0.3em]
\hline
FP = alternativa al Bachillerato & \luecke{4cm} \\[0.3em]
\hline
\end{tabularx}

\vspace{0.5em}
\textbf{Schreibe 3 Vergleichssätze:}

\vspace{0.3em}
1. \es{Mientras que en España} \luecke{9cm}

\vspace{0.4cm}
2. \es{En Alemania hay Gymnasium, ya que} \luecke{7cm}

\vspace{0.4cm}
3. \es{Además,} \luecke{11cm}

\vspace{1em}

\aufgabe{Mediación: FP-Flyer erklären. (Buch S. 103, Aufg. 7)}{10}

\begin{buchverweis}
\textbf{Sprachmittlung:} Ein spanischer Freund möchte wissen, was \glqq{}Formación Profesional\grqq{} (Berufsausbildung) ist. Erkläre den Flyer auf Spanisch -- nicht wörtlich übersetzen, sondern sinngemäß!
\end{buchverweis}

\vspace{0.5em}
\textbf{Deutscher Flyer-Inhalt:}
\begin{itemize}[leftmargin=1.5em, nosep]
    \item Duale Ausbildung: Betrieb + Berufsschule
    \item Dauer: 2--3,5 Jahre
    \item Bezahlung während der Ausbildung
    \item Abschluss: Gesellenbrief / IHK-Prüfung
\end{itemize}

\vspace{0.5em}
\textbf{Deine Erklärung auf Spanisch:}

\vspace{0.3em}
\begin{tabularx}{\textwidth}{@{}|X|@{}}
\hline
\\[0.5em]
\luecke{14cm} \\[0.7em]
\luecke{14cm} \\[0.7em]
\luecke{14cm} \\[0.7em]
\luecke{14cm} \\[0.7em]
\luecke{14cm} \\[0.5em]
\hline
\end{tabularx}

\vspace{0.5em}
\textbf{Hilfreiche Wendungen:}
\begin{itemize}[leftmargin=1.5em, nosep]
    \item \es{En Alemania existe un sistema que se llama...}
    \item \es{Los aprendices trabajan en una empresa y además van a...}
    \item \es{El programa dura aproximadamente...}
    \item \es{Al final reciben un certificado que...}
\end{itemize}

\vspace{1em}

\aufgabe{Wortschatz: Bildung und Beruf. Übersetze.}{6}

\begin{tabularx}{\textwidth}{@{}|X|X|X|X|@{}}
\hline
\rowcolor{hellgrau}
\textbf{Spanisch} & \textbf{Deutsch} & \textbf{Spanisch} & \textbf{Deutsch} \\
\hline
el instituto & \luecke{2.5cm} & la formación & \luecke{2.5cm} \\[0.5em]
\hline
el bachillerato & \luecke{2.5cm} & la empresa & \luecke{2.5cm} \\[0.5em]
\hline
la carrera & \luecke{2.5cm} & el aprendiz & \luecke{2.5cm} \\[0.5em]
\hline
\end{tabularx}

\vspace{1em}
\hrule
\vspace{0.3em}
\begin{flushright}
\textbf{Gesamt: \luecke{1cm} / 30 Punkte}
\end{flushright}

\end{document}
